\clearpage
\begin{thebibliography}{9}
\bibitem{hurwitz}
A. Hurwitz, \emph{Ueber die angen\"aherte Darstellung der Irrationalzahlen
durch rationale Br\"uche}, Mathematische Annalen \textbf{39} (1891),
279--284. \url{https://doi.org/10.1007/BF01206656}.
\bibitem{filaseta}
M. Filaseta, \emph{Math 784: Algebraic Number Theory, Instructor's Notes},
Theorem 25, printed pp.~20--21.
\url{https://people.math.sc.edu/filaseta/gradcourses/TheMath784Notes.pdf}.
\bibitem{bv}
V. Beresnevich and S. Velani, \emph{Ubiquity and a general logarithm law
for geodesics}, arXiv:0707.1225 (2007).
\url{https://arxiv.org/abs/0707.1225}.
\bibitem{bdv}
V. Beresnevich, D. Dickinson and S. Velani,
\emph{Measure theoretic laws for lim sup sets}, arXiv:math/0401118,
version 3 (2004). \url{https://arxiv.org/abs/math/0401118}.
\bibitem{repo}
V. Reshetnikov and contributors, \emph{ProveIt: Concise presentations of
intuitionistic primitive recursive arithmetic}. Repository snapshot
\texttt{34c4c35edc8e648e57587c6faf3a5cfdd7879777}; source inspection
15 November 2026. Exact pinned source links and toolchain versions are
included in \path{sources/source_manifest.json}.
\url{https://github.com/VladimirReshetnikov/ProveIt/tree/main/Logic/PrimitiveRecursiveArithmetic/ConciseIPRA}.
\bibitem{requested}
Requested statement references: \emph{Hurwitz's theorem (number theory)}
and the Khinchin section of \emph{Diophantine approximation}, Wikipedia,
accessed 15 November 2026. The proofs above are self-contained;
these links identify the requested classical formulations.
\url{https://en.wikipedia.org/wiki/Hurwitz\%27s_theorem_(number_theory)};
\url{https://en.wikipedia.org/wiki/Diophantine_approximation}.
\end{thebibliography}
